\section*{Input File Syntax and Structure}
All solver setup and configuration is done through the input file, which is passed to the solver executable. The basic input file structure is given below.

\begin{lstlisting}[language=C++]
// A comment

#include "OtherInputSettings.inp"

Section {
    parameter1 = 6.63e-34
    parameter2 = 137

    SubSection {
        configuration = someOption
        subparameter1 = "filename.csv"
        subparameter2 = 1
    }
}
\end{lstlisting}

The \lstinline[language=C++]!#include! directive on line 3 allows the input file to be split up into multiple files. The directive essentially copy and pastes the contents of the included file at the location of the directive. 

Some essential syntax rules:
\begin{itemize}
    \item Everything is case sensitive.

    \item Scientific notation i.e. \lstinline[language=C++]!5.67e-8! can only be used for parameters which are floating point numbers.

    \item The file structure is completely free-form, meaning braces can be put anywhere, and indenting is ignored. However, there must be a new line after a parameter is specified with an ``\lstinline[language=C++]!=!".

    \item Input file options can be put in any order, as long as they follow the correct tree structure.

    \item White space is ignored, except inside double quotes e.g. \lstinline[language=C++]!"white space not ignored here"!.
\end{itemize}

The rest of this documents gives all the required and optional solver parameters that the input file requires. An example of a complete input file is given at the end of this document.


